\documentclass[a4paper,12pt]{ctexart}
\usepackage{xcolor}
\usepackage[top=1.8cm, bottom=1.8cm, left=1.8cm, right=1.8cm]{geometry}
\usepackage{array}
\usepackage{listings}
\usepackage{hyperref}
\hypersetup{colorlinks=true,linkcolor=blue!70!black}
\linespread{1.35}

\lstset{
  basicstyle=\ttfamily\small,
  keywordstyle=\color{blue},
  commentstyle=\color{gray},
  stringstyle=\color{red},
  showstringspaces=false,
  frame=single,
  numbers=left,
  numberstyle=\tiny,
  breaklines=true
}

\title{\textcolor{blue!70!black}{\Huge\textbf{Python 基础}}}
\author{从 Scratch 积木到真正的代码}
\date{}

\begin{document}
\maketitle
\thispagestyle{empty}

\section*{欢迎来到真正的编程！}

你已经在 Scratch 里学会了编程的核心概念，又通过计算思维学会了"像程序员一样思考"。现在，是时候学习\textbf{第一门真正的编程语言}了。

Python 是目前全球最流行的入门编程语言。它的特点是\textbf{语法简单、读起来像英语}。很多大学生、甚至工作中的程序员都在用 Python。

更重要的是：你在 Scratch 里学的每一个概念，在 Python 中都有对应的写法。你不是从零开始——你只是换了一种表达方式。

\section*{第一课：搭建环境与第一个程序}

\subsection*{安装 Python}

\begin{enumerate}
  \item 打开浏览器，访问 \url{https://www.python.org/downloads/}
  \item 下载最新版本的 Python（点击黄色的 Download 按钮）
  \item 安装时，\textbf{一定记得勾选} "Add Python to PATH"
  \item 安装完成后，打开"命令提示符"（按 Win+R，输入 \texttt{cmd}，回车）
  \item 输入 \texttt{python --version}，如果显示版本号，就安装成功了
\end{enumerate}

\subsection*{第一个程序：Hello World}

打开任意文本编辑器（记事本就可以），输入：

\begin{lstlisting}[language=Python]
print("Hello, World!")
\end{lstlisting}

保存为 \texttt{hello.py}。然后在命令行输入：

\begin{lstlisting}[language=bash]
python hello.py
\end{lstlisting}

屏幕上就会显示：

\begin{lstlisting}
Hello, World!
\end{lstlisting}

\textbf{恭喜！你写了第一个 Python 程序！} \texttt{print()} 是 Python 的"说话"指令——它把括号里的内容显示在屏幕上。

\subsection*{对比 Scratch}

在 Scratch 里，想让小猫说话，你用的是"说 你好！ 2 秒"积木。在 Python 里就是 \texttt{print("你好！")}。本质上是一回事。

\subsection*{练习}

\begin{enumerate}
  \item 修改你的 \texttt{hello.py}，让它输出三行：你的名字、你最喜欢的颜色、你今天的日期。
  \item 故意制造一个错误：把 \texttt{print} 写成 \texttt{pront}，运行看看报错信息是什么样的。然后改回来。
  \item \textbf{[思考]} \texttt{print("Hello" + "World")} 和 \texttt{print("Hello", "World")} 的输出有什么不同？动手试试。
\end{enumerate}

\section*{第二课：变量与数据类型}

\subsection*{变量——贴了标签的盒子}

变量就是一个\textbf{存储数据的盒子}，你给它起个名字，往里面放东西。

\begin{lstlisting}[language=Python]
score = 0          # 变量 score 里存了数字 0
name = "小明"      # 变量 name 里存了文字"小明"
price = 12.5      # 变量 price 里存了小数 12.5
is_game_over = False  # 变量 is_game_over 里存了"假"
\end{lstlisting}

\texttt{=} 不是数学里的"等于"，而是\textbf{赋值}——把右边的值装进左边名字的盒子里。

\subsection*{数据类型}

Python 中有几种基本的数据类型：

\begin{center}
\begin{tabular}{|c|c|c|}
\hline
\textbf{类型} & \textbf{名字} & \textbf{例子} \\ \hline
整数 & \texttt{int} & \texttt{42}, \texttt{-3}, \texttt{0} \\ \hline
浮点数 & \texttt{float} & \texttt{3.14}, \texttt{-0.5}, \texttt{2.0} \\ \hline
字符串 & \texttt{str} & \texttt{"你好"}, \texttt{'Python'} \\ \hline
布尔值 & \texttt{bool} & \texttt{True}, \texttt{False} \\ \hline
\end{tabular}
\end{center}

\textbf{注意：} 字符串要用引号括起来，数字不用。\texttt{42} 是数字，\texttt{"42"} 是文字。

\subsection*{Scratch 对照}

在 Scratch 里，你在红色"变量"积木里创建了"得分"——那就是 Python 中的变量。Scratch 会帮你区分数据类型，Python 也会自动识别，但你自己要知道："1" 和 1 是不同的东西。

\subsection*{练习}

\begin{lstlisting}[language=Python]
# 创建三个变量：你的名字、年龄、身高
my_name = "张三"
my_age = 10
my_height = 1.45

# 打印出来
print(my_name)
print(my_age)
print(my_height)
\end{lstlisting}

运行看看。然后试着把 \texttt{my\_age} 改成文字：\texttt{my\_age = "十岁"}。还能正常运行吗？

\section*{第三课：输入与输出}

\subsection*{让程序和你对话}

\begin{lstlisting}[language=Python]
name = input("你叫什么名字？")
print("你好，" + name + "！")
\end{lstlisting}

\texttt{input()} 会在屏幕上显示括号里的文字，然后等待你输入。你输入的内容会被存到变量里。

\texttt{+} 在字符串之间可以把它们\textbf{拼接}起来。

\subsection*{注意：input 返回的是字符串}

\begin{lstlisting}[language=Python]
age = input("你几岁？")      # age 是字符串，比如 "10"
next_year = age + 1         # 报错！不能把字符串和数字相加
\end{lstlisting}

需要用 \texttt{int()} 把字符串转换成整数：

\begin{lstlisting}[language=Python]
age = int(input("你几岁？"))  # age 现在是整数 10
next_year = age + 1          # 正确！结果是 11
print("明年你" + str(next_year) + "岁了")
\end{lstlisting}

\texttt{str()} 又把数字变回字符串，这样就能用 + 拼接了。

\subsection*{Scratch 对照}

Scratch 里的"询问 你叫什么名字？并等待"积木 + "回答"积木 = Python 中的 \texttt{input()}。

\subsection*{练习}

\begin{enumerate}
  \item 写一个程序：用 \texttt{input()} 问用户出生年份，计算并输出用户年龄（假设今年是 2026 年）。
  \item 扩展上题：如果年龄 $\ge$ 18，输出"你已经成年了"；否则输出"你还未成年"。
  \item \textbf{[调试练习]} 运行以下代码，观察报错信息，然后修复它：
\begin{lstlisting}[language=Python]
age = input("你几岁？")
print("明年你" + age + 1 + "岁了")
\end{lstlisting}
\end{enumerate}

\section*{第四课：条件判断——if 语句}

\subsection*{基本语法}

在 Scratch 里，你用橙色"如果……那么"积木做判断。在 Python 里是这样：

\begin{lstlisting}[language=Python]
score = 85

if score >= 60:
    print("及格了！")
else:
    print("不及格……")
\end{lstlisting}

\textbf{注意事项：}
\begin{itemize}
  \item \texttt{if} 后面的条件\textbf{不用括号}，但要以\textbf{冒号}结尾
  \item 条件下面的代码要\textbf{缩进}（通常按 Tab 或 4 个空格）
  \item Python 用缩进来表示代码属于哪一块（不是花括号 \{\}）
\end{itemize}

\subsection*{多分支：if / elif / else}

\begin{lstlisting}[language=Python]
score = 85

if score >= 90:
    grade = "A"
elif score >= 80:
    grade = "B"
elif score >= 70:
    grade = "C"
elif score >= 60:
    grade = "D"
else:
    grade = "F"

print("你的等级是：" + grade)
\end{lstlisting}

\texttt{elif} = else + if，是"否则如果"的意思。

\subsection*{比较运算符}

\begin{center}
\begin{tabular}{|c|c|c|}
\hline
\textbf{运算符} & \textbf{含义} & \textbf{例子} \\ \hline
\texttt{==} & 等于 & \texttt{score == 100} \\ \hline
\texttt{!=} & 不等于 & \texttt{score != 0} \\ \hline
\texttt{>} & 大于 & \texttt{score > 60} \\ \hline
\texttt{<} & 小于 & \texttt{score < 60} \\ \hline
\texttt{>=} & 大于等于 & \texttt{score >= 60} \\ \hline
\texttt{<=} & 小于等于 & \texttt{score <= 100} \\ \hline
\end{tabular}
\end{center}

\textbf{注意：} \texttt{=} 是赋值，\texttt{==} 才是判断是否相等。这是初学者最容易犯的错误！

\subsection*{练习：猜数字}

结合 \texttt{random} 模块和条件判断，写一个猜数字游戏：

\begin{lstlisting}[language=Python]
import random

secret = random.randint(1, 100)
guess = int(input("猜一个 1-100 之间的数："))

if guess == secret:
    print("太厉害了！一次就猜对！")
elif guess > secret:
    print("大了")
else:
    print("小了")
\end{lstlisting}

（这个版本只能猜一次。后面学了循环，你就能让它一直猜直到猜中。）

\subsection*{练习}

\begin{enumerate}
  \item 写一个程序：输入一个分数，输出对应的等级（90以上A，80-89 B，70-79 C，60-69 D，60以下F）。用 \texttt{if/elif/else} 实现。
  \item 判断以下条件表达式的值（True 或 False）：① \texttt{10 > 5 and 3 < 1} ② \texttt{not (7 == 7)} ③ \texttt{5 != 5 or 2 == 2}
  \item \textbf{[陷阱题]} 以下代码为什么会输出意想不到的结果？改正确。
\begin{lstlisting}[language=Python]
x = 10
if x = 10:
    print("x is 10")
\end{lstlisting}
\end{enumerate}

\section*{第五课：循环——让计算机做重复的事}

\subsection*{for 循环——知道要重复几次}

Scratch 中的"重复 10 次"积木，在 Python 中：

\begin{lstlisting}[language=Python]
for i in range(10):
    print("这是第", i+1, "次")
\end{lstlisting}

\texttt{range(10)} 生成 0 到 9 的数字序列。循环变量 \texttt{i} 依次取每个值。

\subsection*{遍历列表}

\begin{lstlisting}[language=Python]
fruits = ["苹果", "香蕉", "橘子"]
for fruit in fruits:
    print("我喜欢吃" + fruit)
\end{lstlisting}

输出：
\begin{lstlisting}
我喜欢吃苹果
我喜欢吃香蕉
我喜欢吃橘子
\end{lstlisting}

\subsection*{while 循环——不知道要重复几次}

\begin{lstlisting}[language=Python]
secret = 42
guess = 0

while guess != secret:
    guess = int(input("猜数字："))
    
print("猜对了！")
\end{lstlisting}

只要条件成立，循环就一直执行。当 \texttt{guess == secret} 时，循环结束。

\textbf{小心无限循环！} 如果条件永远为真，程序就会永远跑下去。

\subsection*{Scratch 对照}

\begin{center}
\begin{tabular}{|c|c|}
\hline
\textbf{Scratch 积木} & \textbf{Python} \\ \hline
"重复 10 次" & \texttt{for i in range(10):} \\ \hline
"重复执行" & \texttt{while True:} \\ \hline
"重复执行直到……" & \texttt{while 条件:} \\ \hline
\end{tabular}
\end{center}

\subsection*{练习：完善猜数字游戏}

用 \texttt{while} 循环改造刚才的猜数字游戏，让玩家可以一直猜，直到猜中为止。并记录猜了多少次。

\subsection*{练习}

\begin{enumerate}
  \item 用 \texttt{for} 循环打印 1 到 100 之间所有 7 的倍数。
  \item 用 \texttt{while} 循环计算 1 到 100 所有整数之和，并与公式 $n(n+1)/2$ 的结果对比。
  \item \textbf{[代码理解]} 以下循环会输出什么？运行前先推理：
\begin{lstlisting}[language=Python]
for i in range(5):
    if i % 2 == 0:
        print(str(i) + "是偶数")
    else:
        print(str(i) + "是奇数")
\end{lstlisting}
\end{enumerate}

\section*{第六课：列表——装很多东西的盒子}

\subsection*{什么是列表}

列表就是\textbf{有序的多个值的集合}。在 Scratch 里叫"列表"（List），在 Python 里也叫列表。

\begin{lstlisting}[language=Python]
# 创建列表
scores = [95, 87, 92, 78, 88]
names = ["小明", "小红", "小刚"]
mixed = [1, "hello", 3.14, True]  # 不同数据类型也可以

# 访问元素（从 0 开始数！）
print(scores[0])   # 第一个元素：95
print(scores[2])   # 第三个元素：92
print(scores[-1])  # 最后一个元素：88
\end{lstlisting}

\textbf{重要：} 计算机从 0 开始数数，不是从 1。所以第一个元素的索引是 0，第二个是 1……这是初学者最容易搞混的地方。

\subsection*{列表常用操作}

\begin{lstlisting}[language=Python]
nums = [1, 2, 3]

nums.append(4)        # 在末尾添加：nums = [1, 2, 3, 4]
nums.insert(0, 0)     # 在位置 0 插入：nums = [0, 1, 2, 3, 4]
nums.remove(2)        # 删除第一个出现的 2
last = nums.pop()     # 弹出最后一个元素
length = len(nums)    # 获取列表长度
nums.sort()           # 排序
\end{lstlisting}

\subsection*{遍历列表的几种方式}

\begin{lstlisting}[language=Python]
fruits = ["苹果", "香蕉", "橘子"]

# 方式 1：直接遍历元素
for fruit in fruits:
    print(fruit)

# 方式 2：通过索引遍历
for i in range(len(fruits)):
    print(i, fruits[i])

# 方式 3：同时得到索引和元素
for i, fruit in enumerate(fruits):
    print(i, fruit)
\end{lstlisting}

\subsection*{练习：成绩统计}

有一个列表 \texttt{scores = [85, 92, 78, 95, 88, 76, 91]}，请计算：
\begin{itemize}
  \item 总分
  \item 平均分
  \item 最高分
  \item 最低分
\end{itemize}

提示：用 \texttt{sum()}, \texttt{max()}, \texttt{min()} 函数。

\subsection*{练习}

\begin{enumerate}
  \item 创建一个包含 10 个随机整数的列表，输出：最大值、最小值、平均值、中位数。
  \item 有两个列表 \texttt{a = [1, 2, 3, 4, 5]} 和 \texttt{b = [3, 4, 5, 6, 7]}，用列表的方式找出同时在两个列表中的元素（交集）。
  \item \textbf{[挑战]} 写一个程序生成"斐波那契数列"的前 20 项并存到列表中，然后输出。斐波那契：$F_1=1, F_2=1, F_n = F_{n-1} + F_{n-2}$。
\end{enumerate}

\section*{第七课：元组与字典}

\subsection*{元组——不可修改的列表}

元组和列表很像，但创���后\textbf{不能修改}（不能添加、删除、修改元素）。

\begin{lstlisting}[language=Python]
point = (3, 5)        # 一个坐标点
rgb = (255, 128, 0)   # 一种颜色

print(point[0])       # 3
# point[0] = 10       # 错误！元组不能修改
\end{lstlisting}

什么时候用元组？当你确定数据不会改变时——比如坐标、颜色、日期等。

\subsection*{字典——用名字索引}

列表用数字做索引，字典用\textbf{任意名字}做索引。

\begin{lstlisting}[language=Python]
student = {
    "name": "小明",
    "age": 10,
    "score": 95
}

print(student["name"])    # 小明
print(student["score"])   # 95

student["score"] = 98     # 修改
student["class"] = "4班"  # 添加新键
del student["age"]        # 删除
\end{lstlisting}

\textbf{类比：} 字典就像一本真正的字典——你用"词"（键）去查"释义"（值）。

\subsection*{遍历字典}

\begin{lstlisting}[language=Python]
student = {"name": "小明", "age": 10, "score": 95}

for key in student:
    print(key, student[key])

for key, value in student.items():
    print(key, value)
\end{lstlisting}

\subsection*{Scratch 对照}

Scratch 里只有"变量"和"列表"，没有字典。字典是 Python 新增的功能。但在 Scratch 中，你可以用两个列表模拟字典（一个存键、一个存值）。

\subsection*{练习：通讯录}

创建一个字典来存储 3 个同学的姓名和电话号码。然后让用户输入名字，程序显示对应的电话号码。

\subsection*{练习}

\begin{enumerate}
  \item 有一个字典 \texttt{student = \{"name": "小明", "scores": [85, 92, 78]\}}，请计算该学生的平均分并输出。
  \item 用字典统计一个字符串中每个字母出现的次数。例如 \texttt{"hello"} → \texttt{\{'h':1, 'e':1, 'l':2, 'o':1\}}。
  \item \textbf{[对比思考]} 列表和字典有什么不同？什么时候用列表，什么时候用字典？各举一个例子。
\end{enumerate}

\section*{第八课：函数——自制积木}

\subsection*{为什么要用函数？}

你还记得 Scratch 里的"自制积木"吗？你在那里定义了一段代码，然后可以重复使用。Python 里的\textbf{函数}是一样的。

\begin{lstlisting}[language=Python]
def greet(name):
    print("你好，" + name + "！")
    print("欢迎来到 Python 世界！")

# 调用函数
greet("小明")
greet("小红")
\end{lstlisting}

\subsection*{函数结构}

\begin{lstlisting}[language=Python]
def 函数名(参数1, 参数2, ...):
    """函数的说明文档（可选）"""
    # 函数体
    return 返回值  # 可选
\end{lstlisting}

\subsection*{有返回值的函数}

\begin{lstlisting}[language=Python]
def add(a, b):
    return a + b

result = add(3, 5)
print(result)  # 8
\end{lstlisting}

\texttt{return} 会把结果"送回来"给调用者。没有 \texttt{return} 的函数返回 \texttt{None}。

\subsection*{例子：面积计算器}

\begin{lstlisting}[language=Python]
def rectangle_area(width, height):
    return width * height

def circle_area(radius):
    return 3.14 * radius * radius

print("长方形面积：", rectangle_area(5, 3))
print("圆形面积：", circle_area(2))

# 用循环调用函数
for r in range(1, 6):
    print(f"半径{r}的圆面积：{circle_area(r)}")
\end{lstlisting}

\subsection*{Scratch 对照}

\begin{center}
\begin{tabular}{|c|c|}
\hline
\textbf{Scratch} & \textbf{Python} \\ \hline
"自制积木" & \texttt{def 函数名():} \\ \hline
积木的"输入参数" & 函数的参数 \\ \hline
"将……返回"积木 & \texttt{return} 语句 \\ \hline
\end{tabular}
\end{center}

\subsection*{练习：计算器函数}

写一个函数 \texttt{calculator(a, b, op)}，接受两个数字和一个运算符（\texttt{"+"}, \texttt{"-"}, \texttt{"*"}, \texttt{"/"}），返回计算结果。除数为 0 时返回 \texttt{"除数不能为 0"}。

\subsection*{练习}

\begin{enumerate}
  \item 写一个函数 \texttt{is\_palindrome(s)}，判断字符串是否是回文（正反读一样，如 \texttt{"aba"}、\texttt{"上海自来水来自海上"}）。
  \item 写一个函数 \texttt{fib(n)}，返回斐波那契数列的第 $n$ 项。用循环实现，不要用递归。
  \item \textbf{[代码复用]} 你在本课写了多个函数（计算器、回文判断等），现在把它们放在一个文件里，然后写另一个文件来 \texttt{import} 并使用它们。体验模块化的好处。
\end{enumerate}

\section*{第九课：模块——借用别人的代码}

\subsection*{导入模块}

Python 有大量的\textbf{现成模块}供你使用。你不必什么都自己写。

\begin{lstlisting}[language=Python]
import random      # 随机数模块
import math        # 数学模块
import os          # 操作系统模块
\end{lstlisting}

\subsection*{random 模块}

\begin{lstlisting}[language=Python]
import random

print(random.randint(1, 100))     # 1-100 的随机整数
print(random.choice(["石头", "剪刀", "布"]))  # 随机选一个
print(random.shuffle([1, 2, 3, 4, 5]))        # 打乱列表
\end{lstlisting}

\subsection*{math 模块}

\begin{lstlisting}[language=Python]
import math

print(math.sqrt(16))       # 平方根：4.0
print(math.floor(3.7))     # 向下取整：3
print(math.ceil(3.2))      # 向上取整：4
print(math.pi)             # 圆周率：3.14159...
\end{lstlisting}

\subsection*{小练习：石头剪刀布}

\begin{lstlisting}[language=Python]
import random

choices = ["石头", "剪刀", "布"]
computer = random.choice(choices)
player = input("请出拳（石头/剪刀/布）：")

print("电脑出了：" + computer)

if player == computer:
    print("平局！")
elif (player == "石头" and computer == "剪刀") or \
     (player == "剪刀" and computer == "布") or \
     (player == "布" and computer == "石头"):
    print("你赢了！")
else:
    print("你输了！")
\end{lstlisting}

\subsection*{练习}

\begin{enumerate}
  \item 用 \texttt{random} 模块写一个程序：生成一个 4 位数的随机密码（每位是 0-9 的数字）。
  \item 用 \texttt{math} 模块计算半径为 5 的圆的面积和周长（$\pi$ 用 \texttt{math.pi}）。
  \item \textbf{[项目整合]} 把第四课的猜数字游戏和第九课的石头剪刀布放在同一个程序中，让用户可以选择玩哪个游戏。
\end{enumerate}

\section*{第十课：文件读写——让程序记住东西}

现在你的程序一关闭，数据就丢了。\textbf{文件}可以让数据长久保存。

\subsection*{写文件}

\begin{lstlisting}[language=Python]
# 写入文件（覆盖模式）
with open("notes.txt", "w", encoding="utf-8") as f:
    f.write("这是第一行\n")
    f.write("这是第二行\n")

# 追加模式
with open("notes.txt", "a", encoding="utf-8") as f:
    f.write("这是追加的行\n")
\end{lstlisting}

\subsection*{读文件}

\begin{lstlisting}[language=Python]
# 读取全部内容
with open("notes.txt", "r", encoding="utf-8") as f:
    content = f.read()
    print(content)

# 逐行读取
with open("notes.txt", "r", encoding="utf-8") as f:
    for line in f:
        print("行：", line.strip())  # strip() 去掉换行符
\end{lstlisting}

\texttt{with open(...) as f:} 是一个固定写法——打开文件、使用、自动关闭。不用担心忘记关文件。

\subsection*{文件模式}

\begin{center}
\begin{tabular}{|c|c|}
\hline
\textbf{模式} & \textbf{含义} \\ \hline
\texttt{"r"} & 读取（文件必须存在） \\ \hline
\texttt{"w"} & 写入（覆盖已有内容，文件不存在则创建） \\ \hline
\texttt{"a"} & 追加（在文件末尾添加） \\ \hline
\texttt{"r+"} & 读取和写入 \\ \hline
\end{tabular}
\end{center}

\subsection*{小项目：记账本}

\begin{lstlisting}[language=Python]
def add_expense(amount, note):
    with open("expenses.txt", "a", encoding="utf-8") as f:
        f.write(f"{amount},{note}\n")

def show_expenses():
    total = 0
    with open("expenses.txt", "r", encoding="utf-8") as f:
        for line in f:
            amt, note = line.strip().split(",")
            total += int(amt)
            print(f"{note}: {amt}元")
    print(f"总计：{total}元")

# 主程序
while True:
    print("\n1. 添加支出  2. 查看账单  3. 退出")
    choice = input("请选择：")
    if choice == "1":
        amt = int(input("金额："))
        note = input("说明：")
        add_expense(amt, note)
    elif choice == "2":
        show_expenses()
    elif choice == "3":
        break
\end{lstlisting}

\subsection*{练习}

\begin{enumerate}
  \item 写一个程序：让用户输入多行文字（输入"quit"结束），把所有内容保存到 \texttt{notes.txt} 中。
  \item 写一个程序：读取 \texttt{notes.txt} 并统计总行数、总字数和平均每行字数。
  \item \textbf{[扩展]} 给记账本添加一个功能：显示"本月支出最多的类别"。提示：需要记录每个支出的类别（如"餐饮""交通""娱乐"）。
\end{enumerate}

\section*{综合项目：猜数字完整版}

把之前零零散散的知识串起来，做一个完整的猜数字游戏：

\begin{lstlisting}[language=Python]
import random

def guess_game():
    secret = random.randint(1, 100)
    attempts = 0
    best_score = None
    
    # 尝试读取历史最佳记录
    try:
        with open("best_score.txt", "r") as f:
            best_score = int(f.read())
            print(f"历史最佳：{best_score}次猜中")
    except:
        print("还没有历史记录")
    
    print("猜一个 1 到 100 之间的数")
    
    while True:
        guess = int(input("你的猜测："))
        attempts += 1
        
        if guess == secret:
            print(f"猜对了！你用了{attempts}次")
            if best_score is None or attempts < best_score:
                print("新纪录！")
                with open("best_score.txt", "w") as f:
                    f.write(str(attempts))
            break
        elif guess > secret:
            print("大了")
        else:
            print("小了")

guess_game()
\end{lstlisting}

这个游戏包含了：变量、循环、条件判断、函数、文件读写、异常处理——几乎你学过的所有内容！

\section*{总结：你学到了什么？}

\begin{center}
\begin{tabular}{|c|c|}
\hline
\textbf{Scratch 概念} & \textbf{Python 实现} \\ \hline
变量 & \texttt{score = 0} \\ \hline
"如果……那么" & \texttt{if 条件:} \\ \hline
"重复……次" & \texttt{for i in range(n):} \\ \hline
"重复执行" & \texttt{while True:} \\ \hline
"重复执行直到" & \texttt{while 条件:} \\ \hline
列表 & \texttt{[1, 2, 3]} \\ \hline
自制积木 & \texttt{def 函数名():} \\ \hline
"询问并等待" & \texttt{input()} \\ \hline
"说……" & \texttt{print()} \\ \hline
\end{tabular}
\end{center}

\textbf{下一步：} 掌握了 Python 基础之后，你有几个方向可以选择：
\begin{itemize}
  \item \textbf{数据结构与算法} —— 学更多数据结构和经典算法
  \item \textbf{Python → C++} —— 学第二门语言，打开新世界
  \item \textbf{Web 开发} —— 做出自己的网站
  \item \textbf{AI 入门} —— 了解人工智能
\end{itemize}

这些模块都在本系列中！按照《整体课程大纲》的路线继续前进吧。

\section*{\textcolor{blue!70!black}{各课练习参考答案}}

\subsection*{第一课：搭建环境与第一个程序}
\begin{enumerate}
  \item 修改后的 \texttt{hello.py} 应包含三行 \texttt{print()}，分别输出姓名、颜色、日期。
  \item 把 \texttt{print} 写成 \texttt{pront} 会报 \texttt{NameError: name 'pront' is not defined}。
  \item \texttt{print("Hello" + "World")} 输出 \texttt{HelloWorld}（拼接无空格）。\texttt{print("Hello", "World")} 输出 \texttt{Hello World}（默认空格分隔）。
\end{enumerate}

\subsection*{第二课：变量与数据类型}
\begin{enumerate}
  \item 创建三个变量并打印。尝试 \texttt{my\_age = "十岁"} 仍可运行，因为 Python 是动态类型——但之后 age+1 会报错。
  \item 合法变量名：\texttt{score}、\texttt{player\_1}、\texttt{\_temp}。不合法：\texttt{2player}（数字开头）、\texttt{my-name}（含连字符）、\texttt{class}（关键字）。
  \item \texttt{print("10" + 5)} 会报 \texttt{TypeError}：不能拼接字符串和整数。\texttt{print(10 + 5)} 输出 15。
\end{enumerate}

\subsection*{第三课：输入与输出}
\begin{enumerate}
  \item 参考代码：
\begin{lstlisting}[language=Python]
birth_year = int(input("请输入出生年份："))
age = 2026 - birth_year
print("你今年" + str(age) + "岁")
\end{lstlisting}
  \item 在输出前加判断：\texttt{if age >= 18: print("你已经成年了") else: print("你还未成年")}。
  \item \texttt{age} 是字符串，不能直接加整数。修复：\texttt{print("明年你" + str(int(age) + 1) + "岁了")}。
\end{enumerate}

\subsection*{第四课：条件判断}
\begin{enumerate}
  \item 参考代码：
\begin{lstlisting}[language=Python]
score = int(input("输入分数："))
if score >= 90: print("A")
elif score >= 80: print("B")
elif score >= 70: print("C")
elif score >= 60: print("D")
else: print("F")
\end{lstlisting}
  \item ① \texttt{10 > 5 and 3 < 1} → \texttt{True and False} → \texttt{False}。② \texttt{not (7 == 7)} → \texttt{not True} → \texttt{False}。③ \texttt{5 != 5 or 2 == 2} → \texttt{False or True} → \texttt{True}。
  \item 错误：\texttt{if x = 10:} 用了赋值 \texttt{=} 而不是判等 \texttt{==}。修复：\texttt{if x == 10:}。
\end{enumerate}

\subsection*{第五课：循环}
\begin{enumerate}
  \item \texttt{for i in range(1, 101): if i \% 7 == 0: print(i)}。
  \item \texttt{total = 0; i = 1; while i <= 100: total += i; i += 1; print(total)} 输出 5050。
  \item 输出：
\begin{verbatim}
0是偶数
1是奇数
2是偶数
3是奇数
4是偶数
\end{verbatim}
\end{enumerate}

\subsection*{第六课：列表}
\begin{enumerate}
  \item 参考代码：
\begin{lstlisting}[language=Python]
import random
nums = [random.randint(1, 100) for _ in range(10)]
print("最大", max(nums), "最小", min(nums), "平均", sum(nums)/len(nums))
sorted_nums = sorted(nums)
n = len(sorted_nums)
median = sorted_nums[n//2] if n%2==1 else (sorted_nums[n//2-1]+sorted_nums[n//2])/2
\end{lstlisting}
  \item 交集：\texttt{list(set(a) \& set(b))} → [3, 4, 5]。
  \item 参考代码：
\begin{lstlisting}[language=Python]
fib = [1, 1]
for i in range(2, 20):
    fib.append(fib[i-1] + fib[i-2])
print(fib)
\end{lstlisting}
\end{enumerate}

\subsection*{第七课：元组与字典}
\begin{enumerate}
  \item 平均分 = \texttt{sum(student["scores"]) / len(student["scores"])} = (85+92+78)/3 = 85。
  \item 参考代码：
\begin{lstlisting}[language=Python]
s = "hello"
count = {}
for ch in s:
    count[ch] = count.get(ch, 0) + 1
\end{lstlisting}
  \item 列表用数字索引，适合有序集合（如排行榜）。字典用键索引，适合快速查找（如根据名字查电话）。
\end{enumerate}

\subsection*{第八课：函数}
\begin{enumerate}
  \item 参考代码：
\begin{lstlisting}[language=Python]
def is_palindrome(s):
    return s == s[::-1]
\end{lstlisting}
  \item 参考代码：
\begin{lstlisting}[language=Python]
def fib(n):
    a, b = 1, 1
    for _ in range(n-1):
        a, b = b, a + b
    return a
\end{lstlisting}
  \item 将函数定义放在 \texttt{utils.py} 中，在 \texttt{main.py} 中用 \texttt{from utils import calculator, is\_palindrome} 导入使用。
\end{enumerate}

\subsection*{第九课：模块}
\begin{enumerate}
  \item \texttt{import random; password = "".join(str(random.randint(0, 9)) for _ in range(4))}。
  \item \texttt{import math; area = math.pi * 5**2; circumference = 2 * math.pi * 5}。
  \item 用一个主菜单，\texttt{input()} 获取用户选择，用 \texttt{if/elif} 分支到不同的游戏。
\end{enumerate}

\subsection*{第十课：文件读写}
\begin{enumerate}
  \item 参考代码：
\begin{lstlisting}[language=Python]
with open("notes.txt", "w", encoding="utf-8") as f:
    while True:
        line = input()
        if line == "quit": break
        f.write(line + "\n")
\end{lstlisting}
  \item 参考代码：
\begin{lstlisting}[language=Python]
with open("notes.txt", "r", encoding="utf-8") as f:
    lines = f.readlines()
    total_chars = sum(len(line.strip()) for line in lines)
    print(f"行数：{len(lines)}，总字数：{total_chars}，平均：{total_chars/len(lines):.1f}")
\end{lstlisting}
  \item 添加一个"类别"字段到记录中，然后用字典统计每类支出总和。
\end{enumerate}

\section*{综合练习与答案}

\subsection*{基础题}

\begin{enumerate}
  \item 写一个程序，输入三个数，输出其中最大的数。
  \item 用循环打印 1 到 100 之间所有能被 3 整除但不能被 5 整除的数。
  \item 写一个函数 \texttt{is\_palindrome(s)}，判断字符串是否是回文（正着读和反着读一样，如 "aba"、"上海自来水来自海上"）。
  \item 生成一个包含 20 个随机数的列表，输出其中的奇数和偶数分别有多少个。
  \item 读一个文本文件，统计其中有多少个单词（用 \texttt{split()} 即可）。
\end{enumerate}

\subsection*{参考答案}

\textbf{练习 1：}
\begin{lstlisting}[language=Python]
a = int(input("第一个数："))
b = int(input("第二个数："))
c = int(input("第三个数："))
print(max(a, b, c))
\end{lstlisting}

\textbf{练习 2：}
\begin{lstlisting}[language=Python]
for i in range(1, 101):
    if i % 3 == 0 and i % 5 != 0:
        print(i, end=" ")
\end{lstlisting}

\textbf{练习 3：}
\begin{lstlisting}[language=Python]
def is_palindrome(s):
    return s == s[::-1]

print(is_palindrome("aba"))   # True
print(is_palindrome("hello")) # False
\end{lstlisting}

\textbf{练习 4：}
\begin{lstlisting}[language=Python]
import random
nums = [random.randint(1, 100) for _ in range(20)]
odds = [n for n in nums if n % 2 == 1]
evens = [n for n in nums if n % 2 == 0]
print(f"奇数：{len(odds)}个, 偶数：{len(evens)}个")
\end{lstlisting}

\textbf{练习 5：}
\begin{lstlisting}[language=Python]
with open("article.txt", "r", encoding="utf-8") as f:
    text = f.read()
words = text.split()
print(f"单词数：{len(words)}")
\end{lstlisting}

\section*{扩展项目}

\subsection*{项目一：学生成绩管理系统}

要求：用字典存储学生信息（学号、姓名、各科成绩），实现以下功能：
\begin{enumerate}
  \item 添加学生
  \item 删除学生
  \item 修改成绩
  \item 按总分排名
  \item 按科目统计平均分
  \item 数据保存到文件（下次启动时自动加载）
\end{enumerate}
用函数组织代码，每个功能一个函数。

\textbf{关联模块：} 本项目用到文件 IO（本模块第十课）和字典（本模块第七课）。完成后可以对比 \texttt{ds-algo/} 中的排序算法，看你用的排序方法和课程中的哪种算法一致。

\subsection*{项目二：密码生成器}

要求：
\begin{enumerate}
  \item 用户指定密码长度和包含的字符类型（大小写字母/数字/符号）
  \item 生成随机密码
  \item 评估密码强度（弱/中/强）
\end{enumerate}

\textbf{关联模块：} 本项目用到 \texttt{random} 模块（本模块第九课）。密码强度的数学原理见 \texttt{computer\_maths/} 的组合数学章节。

\section*{本模块与其他模块的关联}

\begin{itemize}
  \item \textbf{计算思维} → 本模块：你已经在计算思维模块学会了画流程图和写伪代码。现在你把这些伪代码变成了真正的 Python 程序。回头看你在计算思维模块写的伪代码，试着用 Python 实现它。
  \item \textbf{计算机数学} → 本模块：Python 中的布尔运算（\texttt{and/or/not}）对应 \texttt{computer\_maths/} 中的布尔代数；Python 的 \texttt{set} 类型对应集合论。学完 \texttt{computer\_maths/} 再来写 Python 代码，你会理解 "为什么这样设计"。
  \item 本模块 → \textbf{数据结构与算法}：本模块的列表、字典、函数是 ds-algo 模块的基础。你在本模块写的猜数字游戏用了线性搜索，ds-algo 模块会教你二分搜索——快了 100 倍。
  \item 本模块 → \textbf{Web 开发}：本模块的 Flask 入门需要 Python 基础。你在这里学会的 Python 语法会被直接用到 Web 后端开发中。
  \item 本模块 → \textbf{AI 入门}：AI 模块的代码全部用 Python 实现。你的 Python 基础决定了你能多深入理解 AI 的内部实现。
\end{itemize}

\end{document}
